\section{Limitations and Conclusion}
\label{sec:limitations}

The current evidence is strongest for controlled local-SGD ledgers and small
adaptive shift settings. The language-model experiments are tiny, and the
multi-domain suite is synthetic. These experiments are useful because they
isolate the mechanism and its cost accounting, but they do not replace
large-scale real-data runs. The paper should therefore claim a mechanism and a
fixed-compute proof of concept, not a universal recipe for foundation-scale
training.

Probe design is also central. If the probe is too noisy, too expensive, or
misaligned with the target, the controller can overfit the probe or lose to a
simple external sweep. E096, E089, and E090 are included for exactly this
reason. They mark regimes where probe labels are not enough, external HPO is
too strong, or the stale-default comparator is too severe.

Finally, deliberate diversity has to be real. Under exchangeable IID workers,
the theory and null audit both say that arm labels cannot carry persistent
hyperparameter information. The useful design space is the one where arms
encode different optimizer moves and where the probe can cheaply identify
which move is useful for the present shared trajectory.

In summary, distributed optimization already exposes parallel candidate
updates. Deliberately diverse aggregation turns those candidates into an
explore-exploit mechanism: create meaningful optimizer-arm diversity, score
candidate next states with cheap probes, and reinvest the result into one
model. The fixed-compute outer-control experiments show that this can beat
external HPO, and the adaptive shift experiments show the same idea working
across language-model and multi-domain task families.
